\begin{abstract}
Prediction Shift Backdoor Detection (PSBD) flags a poisoned input by how little
its prediction moves when the network is perturbed at inference. It was built for
ResNets, where the residual stream is the one obvious place to put that
perturbation, and a transformer block offers many places and several kinds of
perturbation. We show that this choice, rather than the method, decides whether
the detector works. Masking whole tokens at the input of every attention block
reaches a mean AUROC of \HeadlineAurocAdaptive over \PanelCellsCached backdoored
ViT-B/16 models spanning \PanelDatasetsCached datasets and poison rates from 1\% to
10\%, against \PublishedAurocAdaptive for the placement the original work used.
The paired gain is \HeadlineGainAdaptiveAuroc
[\HeadlineGainAdaptiveAurocLow, \HeadlineGainAdaptiveAurocHigh] and it grows to
\HardGainAdaptiveAuroc on the attacks that resist detection. Where the
perturbation goes and what it injects are separate axes, and what decides the
operator is the side of the LayerNorm the perturbation lands on. Injected before
a LayerNorm, Gaussian noise costs \GaussianMinusTokenMaskAttentionNorm AUROC
against token masking at the attention input. Injected after a LayerNorm, the 2 operators
tie. We explain that by locating the backdoor, which is 1 direction in the
residual stream, written in the last third of the network and routed through
attention from the trigger's own tokens to the class token. Removing whole tokens
before attention removes that route, while noise at the same site is
renormalized away by the LayerNorm behind it. We also report the limits. An attacker who trains against 1 probe
defeats it, a union of probes restores detection and all-to-all attacks remain
open.
\end{abstract}
